\documentclass[10pt]{article}

\usepackage[a4paper,left=23mm,right=23mm,top=20mm,bottom=22mm]{geometry}
\usepackage{amsmath,amssymb,amsthm,mathtools}
\usepackage{microtype}
\usepackage[hidelinks]{hyperref}

\newtheorem{theorem}{Theorem}
\newcommand{\Z}{\mathbb Z}
\newcommand{\Qform}{\mathcal Q}

\setlength{\parindent}{1.3em}
\setlength{\parskip}{0pt}
\emergencystretch=2em

\hypersetup{
  pdftitle={An Infinite Pell Family for z2 + xy2 + x3 + 3 = 0},
  pdfauthor={Jamal Agbanwa}
}

\begin{document}

\begin{center}
{\LARGE\bfseries An Infinite Pell Family for $z^2+xy^2+x^3+3=0$\par}
\vspace{0.55em}
\end{center}
\vspace{0.5em}

\begin{abstract}
We prove, by an explicit elementary construction, that
\[
 z^2+xy^2+x^3+3=0
\]
has infinitely many integer solutions.  Fixing $x=-3$ reduces the equation to the Pell-type conic $z^2-3y^2=24$.  A norm-preserving recurrence, obtained from the unit $2+\sqrt3$, then generates an infinite strictly increasing sequence of solutions.  We also explain in detail how this useful slice and its recurrence arise.
\end{abstract}

\section{The explicit infinite family}

\begin{theorem}\label{thm:main}
Define integer sequences $(y_n)_{n\ge 0}$ and $(z_n)_{n\ge 0}$ by
\begin{equation}\label{eq:recurrence}
 y_0=2,\qquad z_0=6,
 \qquad
 y_{n+1}=2y_n+z_n,\qquad
 z_{n+1}=3y_n+2z_n.
\end{equation}
Then, for every $n\ge 0$ and every $\varepsilon,\delta\in\{\pm1\}$,
\begin{equation}\label{eq:family}
 (x,y,z)=\bigl(-3,\varepsilon y_n,\delta z_n\bigr)
\end{equation}
is an integer solution of
\begin{equation}\label{eq:surface}
 z^2+xy^2+x^3+3=0.
\end{equation}
Moreover, the solutions $(-3,y_n,z_n)$ are pairwise distinct.  Consequently, \eqref{eq:surface} has infinitely many integer solutions.
\end{theorem}

\begin{proof}
Put
\[
 \Qform(y,z)=z^2-3y^2.
\]
A direct calculation gives the invariant identity
\begin{align*}
 \Qform(2y+z,3y+2z)
 &=(3y+2z)^2-3(2y+z)^2\\
 &=z^2-3y^2
 =\Qform(y,z).
\end{align*}
Since
\[
 \Qform(y_0,z_0)=6^2-3\cdot2^2=24,
\]
induction using \eqref{eq:recurrence} yields
\begin{equation}\label{eq:pell-invariant}
 z_n^2-3y_n^2=24
 \qquad(n\ge0).
\end{equation}
On taking $x=-3$, the left-hand side of \eqref{eq:surface} becomes
\[
 z_n^2-3y_n^2-27+3
 =24-24=0
\]
after applying \eqref{eq:pell-invariant}.  Because only $y^2$ and $z^2$ occur in \eqref{eq:surface}, the two signs in \eqref{eq:family} may be chosen independently.

It remains only to verify infinitude.  The initial values are positive, and \eqref{eq:recurrence} therefore implies inductively that $y_n,z_n>0$ for every $n$.  Hence
\[
 y_{n+1}=2y_n+z_n>y_n,
\]
so $(y_n)$ is strictly increasing.  Thus the triples $(-3,y_n,z_n)$ are pairwise distinct, proving the assertion.
\end{proof}

The first positive solutions furnished by the theorem are
\[
 (-3,2,6),\quad(-3,10,18),\quad(-3,38,66),\quad
 (-3,142,246),\quad(-3,530,918),\ldots.
\]

\section{How the construction was found}\label{sec:discovery}

The decisive step is to expose an indefinite quadratic form hidden inside the cubic equation.  First, no solution can have $x\ge0$: in that case $z^2$, $xy^2$, and $x^3$ are all nonnegative, so the left-hand side of \eqref{eq:surface} is at least $3$.  Every integer solution must therefore have $x<0$.

Write $x=-a$ with $a\in\Z_{>0}$.  Equation \eqref{eq:surface} is then exactly equivalent to
\begin{equation}\label{eq:slices}
 z^2-ay^2=a^3-3.
\end{equation}
For each fixed $a$, the original cubic surface has been cut by the plane $x=-a$ into a Pell-type conic.  The problem is now to choose $a$ so that this conic has one integer point and also admits a nontrivial norm-preserving motion.

The first two choices fail for elementary congruence reasons.  For $a=1$, equation \eqref{eq:slices} reads
\[
 z^2-y^2=-2,
\]
which is impossible modulo $4$, since a difference of two squares is never congruent to $2\pmod4$.  For $a=2$, it reads
\[
 z^2-2y^2=5.
\]
If $y$ is even, this would give $z^2\equiv5\pmod8$; if $y$ is odd, it would give $z^2\equiv7\pmod8$.  Both are impossible because a square modulo $8$ is $0$, $1$, or $4$.

The next choice, $a=3$, succeeds immediately:
\begin{equation}\label{eq:pell}
 z^2-3y^2=24,
\end{equation}
and the small identity
\[
 6^2-3\cdot2^2=24
\]
provides the seed $(y,z)=(2,6)$.  Thus the remaining task is not to find more points independently, but to find a transformation preserving the left-hand side of \eqref{eq:pell}.

The general multiplicative identity
\begin{equation}\label{eq:norm-identity}
 (rz+asy)^2-a(sz+ry)^2
 =(r^2-as^2)(z^2-ay^2)
\end{equation}
is verified by expanding both sides: the mixed terms cancel, leaving
\((r^2-as^2)z^2-a(r^2-as^2)y^2\).  It says that whenever
$r^2-as^2=1$, the map
\[
 (y,z)\longmapsto (sz+ry,\,rz+asy)
\]
preserves $z^2-ay^2$.  For $a=3$, the smallest positive nontrivial solution of
\[
 r^2-3s^2=1
\]
is $(r,s)=(2,1)$.  Substitution into the preceding map gives precisely
\[
 (y,z)\longmapsto(2y+z,\,3y+2z),
\]
which is the recurrence in \eqref{eq:recurrence}.  In the suggestive notation of quadratic norms, this is simply multiplication by $2+\sqrt3$:
\[
 z_{n+1}+y_{n+1}\sqrt3
 =(z_n+y_n\sqrt3)(2+\sqrt3).
\]
Since $(2+\sqrt3)(2-\sqrt3)=1$, this multiplication preserves the norm
$z^2-3y^2$; the elementary calculation in the proof of Theorem~\ref{thm:main} is exactly this norm argument written without invoking any algebraic-number-theoretic machinery.

Iterating gives the compact closed form
\begin{equation}\label{eq:closed-form}
 z_n+y_n\sqrt3=(6+2\sqrt3)(2+\sqrt3)^n,
\end{equation}
which follows at once by induction from the recurrence.  Formula \eqref{eq:closed-form} explains both the rapid growth and the inexhaustibility of the family, while Theorem~\ref{thm:main} supplies the completely elementary proof that every generated triple is integral, solves the original equation, and is distinct from all earlier ones.

\end{document}
